\chapter{Les guides de Chult}

\sessioninfo{Session 6}{Mardi 22 septembre 2026}

En quittant la demeure de Wakanga O'Tamu, les quatre aventuriers possédaient enfin une direction générale. Quelque part dans la jungle se trouvait Vorn, le Gardien Bouclier qu'ils étaient venus chercher. Ailleurs se cachait peut-être le \textit{Soulmonger}, l'artefact responsable du mal qui frappait les ressuscités. Entre eux et ces deux objectifs s'étendait cependant l'essentiel de Chult, dont leur carte ne révélait encore qu'une faible partie.

Avant de s'y aventurer, ils devaient commencer par une tâche beaucoup plus simple : trouver où dormir.

Ils reprirent donc le chemin du quartier du marché. À proximité de la demeure de Wakanga, un petit temple consacré à Gond attira leur attention. Morrÿs associait ce dieu aux serruriers et aux portes, aussi bien celles que l'on franchissait que celles qui s'ouvraient sur le destin. Destos reconnut pour sa part le symbole d'une forge : Gond protégeait également les artisans et tous ceux qui travaillaient de leurs mains.

Le groupe poursuivit sa route sans entrer dans le temple. Au-delà, Port Nyanzaru semblait vivre à deux rythmes. Des rues silencieuses demeuraient suspendues au passage de la procession funèbre ; des habitants vêtus de blanc et de noir rejoignaient le deuil ou l'observaient à distance. Ailleurs, les étoffes éclatantes, les voix fortes et les sonorités du créole local, faites de claquements de gorge et de dents, rendaient à la cité toute son agitation.

Un palais doré dominait les hauteurs, visible depuis les quartiers marchands comme depuis le port. Plus bas, le Grand Souk grondait déjà du bruit des négociations. Kael aurait volontiers plongé dans ses allées, mais Morrÿs ramena le groupe à une priorité plus raisonnable : localiser l'auberge, déposer leurs affaires, puis revenir visiter la ville sans craindre de chercher un toit à la nuit tombée.

Ils profitèrent de cette halte pour répartir les objets importants. Conserver ensemble la carte de Chult et le journal nain mentionnant le constructeur du Gardien Bouclier aurait permis au premier voleur habile de repartir avec toutes leurs pistes. Destos glissa donc la carte dans les poches intérieures de sa robe, où une main indiscrète aurait peu de chances de passer inaperçue.

Le journal revint à Obélik. Le goliath montra une petite bourse solidement attachée contre sa jambe, coupant court aux suggestions de cachettes beaucoup moins élégantes qui commençaient déjà à circuler.

\begin{narrateur}

De son point de vue, tous les problèmes étaient petits.

\end{narrateur}

Un passant leur confirma que les marchés restaient ouverts des premières heures du matin jusque tard dans la soirée. Il leur indiqua également la maison de repos de Kaya, plus à l'est, près du bazar consacré aux achats quotidiens. Son regard s'attarda tout de même sur l'orc, le goliath, le nain et le clerc : même à Port Nyanzaru, un tel groupe ne passait pas tout à fait inaperçu.

À mesure qu'ils avançaient, les rues se vidaient sous l'effet de la procession. Puis l'activité reparut autour d'un second marché, où se mêlaient les odeurs de fruits mûrs, d'épices et de viande de dinosaure. Les étals débordaient de nourriture et de produits usuels, tandis que les Chultais circulaient dans des vêtements légers aux couleurs éclatantes.

Une première auberge se dressait à proximité : le Lézard Tonitruant. Le groupe ne s'y arrêta pas. Quelques rues plus loin, une voix sèche jaillit d'un autre établissement et rappela à deux clients que, s'ils ne payaient pas, ils pouvaient sortir. Les malheureux franchirent la porte la mine basse.

Kael n'eut pas besoin de lire l'enseigne pour comprendre qu'ils venaient de trouver Kaya.

\scenebreak

La maison de repos était simple, propre et bien tenue. Un comptoir occupait la gauche de la salle principale, avec plusieurs boissons locales alignées derrière. Quatre ou cinq tables meublaient le reste de la pièce et, à cette heure située entre le petit déjeuner et le déjeuner, seuls quelques clients terminaient encore leur repas.

Kaya reconnut aussitôt les voyageurs annoncés par Wakanga. Elle leur remit les clés des chambres 40 et 42, toutes deux situées au fond d'un long couloir. Leur hébergement et le repas du soir étaient pris en charge ; toute consommation supplémentaire resterait à leurs frais.

Kael voulut s'assurer qu'ils pouvaient laisser leurs affaires sans risque. La question heurta immédiatement la fierté de la tenancière : il n'y avait aucun problème dans son auberge, et quiconque tenterait d'y voler ou d'y troubler la tranquillité aurait rapidement affaire à la justice de Port Nyanzaru. L'établissement n'avait rien de luxueux, mais tout y respirait l'ordre. Kaya semblait surtout veiller sur sa réputation avec la même attention que sur sa caisse.

Les chambres étaient presque identiques. Deux lits simples encadraient une table de nuit commune ; une armoire faisait face aux couchages et une fenêtre donnait sur une cour intérieure calme. Après les semaines passées à bord de l'\textit{Effleureuse}, cette modestie avait quelque chose de confortable.

Obélik ne se sépara pas de sa hache. Les autres conservèrent également l'essentiel de leur équipement, ne déposant que ce qui les encombrerait inutilement pendant quelques heures en ville. En revenant dans la salle, Kael surprit les murmures de clients qui parlaient des morts et de l'impossibilité de les ramener à la vie.

Plutôt que de les interroger et de risquer de révéler ce qu'ils savaient déjà, le groupe choisit une méthode plus discrète : s'installer à proximité et commander à manger. Pour une pièce d'argent par personne, Kaya proposait un assortiment qui rappelait qu'ils avaient quitté la Côte des Épées : salade de fruits, viande de vélociraptor grillée comme du bacon, omelette rosée, œufs d'autruche, fromage de dimétrodon et salade d'orchidées vinaigrée.

Le repas couvrit convenablement leur curiosité. À la table voisine, les conversations évoquaient la troisième procession funèbre de la semaine et les fortunes dépensées pour des morts qui ne pouvaient désormais plus revenir. Les clients n'en savaient manifestement pas davantage, mais leurs inquiétudes confirmaient que le mal dont souffrait Syndra n'était ni secret ni isolé. Toute la ville commençait déjà à en ressentir les conséquences.

Kaya se détendit en constatant que ses nouveaux pensionnaires consommaient sans chercher querelle. Lorsque vint le moment de régler, Kael préféra payer immédiatement les quatre repas avec une pièce d'or. La tenancière lui rendit six pièces d'argent, désormais mieux disposée à répondre à leurs questions.

Le groupe lui apprit alors qu'il préparait une expédition dans la jungle. Kaya les invita à se rasseoir.

\scenebreak

Son premier conseil n'avait rien de rassurant. Hors de Port Nyanzaru, ils devaient davantage se considérer comme des proies que comme des prédateurs. Des explorateurs avaient essayé de cartographier la jungle, de la sécuriser ou d'y établir des positions durables. Beaucoup avaient échoué ; certains n'étaient jamais revenus.

Les dangers ne manquaient pas : dinosaures, maladies, moustiques et morts-vivants attendaient les voyageurs assez téméraires pour quitter les murailles. Kaya connaissait également Fort Beluarian, au nord, ainsi que plusieurs factions tentant de maintenir des positions sur l'île. Elle refusa toutefois de prétendre connaître une jungle dans laquelle elle ne s'était presque jamais aventurée elle-même.

Les aventuriers envisageaient de suivre la Soshenstar vers le sud. Kaya approuva l'idée d'utiliser un canoë ou une pirogue : la rivière offrirait au moins un repère impossible à perdre entre les arbres et leur permettrait de revenir sur leurs pas. Elle les prévint néanmoins que l'eau ne les protégerait ni des créatures qui y vivaient ni de celles qui rôdaient sur ses rives.

Se perdre restait, à ses yeux, le danger le plus probable. Pour l'éviter, sa recommandation fut immédiate : engager un guide. Tous les explorateurs revenus vivants dont elle avait entendu parler en avaient pris un. Elle ne connaissait personne qui soit revenu sans.

Le message avait le mérite d'être clair.

Les guides dépendaient du prince marchand Jobal et laissaient leurs annonces sur un panneau proche de l'auberge. Kaya ne leur recommanda aucun nom. Elle ajouta simplement que, s'ils partaient en expédition, ils devaient emporter leurs affaires et rendre leurs clés. Après quatorze jours sans nouvelles, elle considérerait les chambres abandonnées et vendrait ce qu'ils y auraient laissé.

Le groupe n'avait encore choisi ni guide ni itinéraire précis, mais une certitude venait de s'imposer : il ne s'enfoncerait pas seul dans la jungle.

% \illustration[0.9\textwidth]{chult_session_6.png}
% {\textit{La carte de Chult : Port Nyanzaru, Vorn et l'itinéraire envisagé le long de la Soshenstar.}}
% {carte_chult_session_6}

\scenebreak

À l'intérieur de la maison de repos, juste à côté de l'entrée, un grand panneau de bois rassemblait les annonces des guides. Plusieurs cartes y étaient suspendues à des anneaux et portaient les noms d'Azaka, d'Eku, de Faroul et Gondolo, de Hew Hackinstone, de Musharib et de Salida. Une dernière annonce, signée River Mist et Flask of Wine, avait manifestement été ajoutée après les autres en les repoussant pour se faire une place.

\illustration[0.95\textwidth]{2026-09-22/panneau-guides-nyanzaru.png}
{\textit{Le panneau des guides découvert à la sortie de la maison de repos de Kaya.}}
{panneau_guides_nyanzaru}

Les formules inscrites au dos des cartes donnaient une première impression de chaque candidat. Azaka prétendait que la jungle la reconnaissait à son odeur plutôt qu'à son nom et proposait de troquer ses honoraires contre une dette. Eku invitait les voyageurs perdus à venir la trouver. Faroul et Gondolo promettaient deux guides pour le prix d'un, accompagnés de leur tricératops. Hew misait sur la bravoure et la loyauté naines, tandis que Musharib attendait son heure pour reconquérir une forge perdue. Salida annonçait une progression rapide, toujours vers le sud et sans détours.

La dernière carte différait des autres jusque dans son écriture. Elle promettait de conduire les voyageurs où ils le souhaitaient, en toute discrétion, sous la signature de River Mist et Flask of Wine. Kael et Destos remarquèrent qu'elle avait été glissée là presque de force. Sans être nécessairement mensongère, elle n'avait rien d'une annonce officielle.

Aucun tarif n'était indiqué. Le groupe préférait également rencontrer les guides avant de confier sa survie à l'un d'eux. Puisque Jobal était un prince marchand, aller frapper directement à sa porte pour de simples renseignements leur parut excessif. Ils décidèrent plutôt de revenir chez Wakanga, dont le domestique saurait peut-être les orienter.

À l'entrée de la demeure, ils retrouvèrent l'homme qui les avait accueillis le matin. Il leur évita une visite inutile chez Jobal : les guides pouvaient être rencontrés directement, et le bureau d'information du port disposait de recommandations pour les voyageurs. Surtout, quelques prospectus plus détaillés se trouvaient encore dans la maison.

Le domestique leur révéla aussi le tarif officiel. Chaque guide demandait cinq pièces d'or par jour et exigeait trente jours de paiement avant le départ.

Cent cinquante pièces d'or.

À eux quatre, les aventuriers étaient arrivés à Chult avec deux cents pièces. Il leur restait pourtant à acheter leur équipement et à trouver une embarcation. Le choix d'un guide ne dépendrait donc pas seulement de ses compétences : il leur faudrait peut-être accepter l'une des missions proposées en échange d'une réduction ou d'une suppression des honoraires.

Le domestique revint avec sept prospectus. Le groupe les étala et entreprit de comparer les promesses de chacun.

\section*{Les guides de Port Nyanzaru}

\subsection*{Azaka Stormfang}

\illustration[0.85\textwidth]{2026-09-22/prospectus/Azaka Stormfang.png}
{\textit{Prospectus d'Azaka Stormfang.}}
{prospectus_azaka}

\begin{quote}\itshape
Si vous avez besoin d'un guide pour vous faire traverser la jungle en vie, ne cherchez pas plus loin. Je suis née dans la jungle. J'en connais bien les dangers et je sais comment les éviter au mieux. Si vous me choisissez, vous pouvez être assurés que votre groupe sera entre de bonnes mains.

Je demande cinq pièces d'or par jour, avec un paiement de trente jours à l'avance.

Pour les aventuriers qui m'aideront dans une quête personnelle, je renoncerai à mes honoraires et les guiderai vers des endroits que personne d'autre n'a découverts. Cette quête nous conduira à Firefinger. Le maléfique ptéranodon qui niche au sommet de Firefinger m'a volé un masque en bois, et je veux le récupérer.
\end{quote}

Azaka offrait donc une solution au problème financier du groupe, mais exigeait en retour une dette dont le danger restait difficile à estimer. Firefinger et le ptéranodon qui l'occupait n'avaient rien d'une simple course à travers la ville.

\subsection*{Eku}

\illustration[0.85\textwidth]{2026-09-22/prospectus/Eku.png}
{\textit{Prospectus d'Eku.}}
{prospectus_eku}

\begin{quote}\itshape
Bien que la perspective de piller les antiquités de Chult puisse être tentante, je m'adresse à ceux qui trouvent leur accomplissement en faisant réellement le bien dans ce monde. Je cherche à débarrasser la jungle du mal. De nombreuses ruines et d'anciens temples doivent être purifiés, afin que les esprits de la nature puissent retourner dans leurs lieux sacrés et y vivre en paix.

Il existe de mauvais imposteurs qui chercheront à vous séduire avec des promesses d'or, mais qui ne feront que conduire vos gorges sous leurs lames !

Mon tarif est le tarif standard : cinq pièces d'or par jour, avec un paiement de trente jours à l'avance. Une petite partie sert à rembourser le prince marchand Jobal. Je donne le reste à des œuvres de charité.
\end{quote}

Le discours d'Eku paraissait sincère, mais ne réglait en rien la question du prix. Son insistance sur les imposteurs rappelait également que les annonces ne suffiraient pas à déterminer à qui faire confiance.

\subsection*{Faroul et Gondolo}

\illustration[0.85\textwidth]{2026-09-22/prospectus/Faroul et Gondolo.png}
{\textit{Prospectus de Faroul et Gondolo.}}
{prospectus_faroul_gondolo}

\begin{quote}\itshape
Pourquoi payer pour un seul guide quand vous pouvez en avoir DEUX pour le même prix ! Nous avons l'intelligence, le matériel et notre propre tricératops !

Les trésors abondent ! Lorsque nous aurons fini de vous conduire vers tous les sites anciens et lieux sacrés que nous connaissons, vous nagerez dans l'or et la magie. Tous les autres guides vous feront tourner en rond et vous laisseront faire tout le travail. Ne vous laissez pas berner !

Lorsque nous retournerons à Port Nyanzaru, vous aurez assez d'histoires et de richesses pour toute une vie.

Notre tarif est de cinq pièces d'or par jour, avec un paiement de trente jours à l'avance, mais nous sommes disposés à réduire notre tarif en échange d'une part égale de tous les trésors découverts.
\end{quote}

La perspective de disposer de deux guides pour le même tarif avait de quoi séduire, tout comme le tricératops. L'animal s'accordait toutefois assez mal avec leur projet de descendre la rivière en canoë, et la promesse de richesses paraissait presque trop enthousiaste.

\subsection*{Hew Hackinstone}

\illustration[0.85\textwidth]{2026-09-22/prospectus/Hew Hackinstone.png}
{\textit{Prospectus de Hew Hackinstone.}}
{prospectus_hew}

\begin{quote}\itshape
J'ai parcouru toutes les montagnes de Chult. Je connais chaque gorge, chaque crevasse, chaque plateau et chaque sentier. Ma mémoire est faite d'acier : je n'oublie rien. Faites confiance à un nain. Comme je le dis toujours, un nain ne vous laissera jamais tomber.

Je suis sans peur. Il n'existe aucun dragon, aucune bête ni aucun mort-vivant capable de m'effrayer. Voulez-vous d'un guide qui s'enfuira au premier signe de danger ? Non. Vous voulez un nain solide, et je suis le plus solide de Chult.

Un dragon rouge m'a arraché un bras et je suis toujours là. Il en faudra bien davantage pour arrêter Hew Hackinstone.

Les cinq pièces d'or par jour que vous dépenserez pour mes services seront parmi les meilleures pièces que vous aurez jamais dépensées.
\end{quote}

La loyauté naine de Hew trouva naturellement un écho chez Morrÿs. Kael se demanda néanmoins s'il ne serait pas préférable d'engager quelqu'un né sur l'île, à moins que le guide ne vive à Chult depuis suffisamment longtemps pour en connaître chaque danger.

\subsection*{Musharib}

\illustration[0.85\textwidth]{2026-09-22/prospectus/Musharib.png}
{\textit{Prospectus de Musharib.}}
{prospectus_musharib}

\begin{quote}\itshape
La jungle est ma maison. Je la connais bien. Je peux vous guider n'importe où dans Chult pour cinq pièces d'or par jour.

Si c'est l'aventure que vous recherchez, rejoignez-moi dans ma quête pour reprendre Hrakhamar, la forge ancestrale de mon clan. Une éruption volcanique nous a forcés à l'abandonner il y a de nombreuses années. Les salamandres de feu qui s'y sont installées depuis doivent être chassées ou détruites.

À tout le moins, nous devons nous infiltrer et récupérer une relique appelée le Gantelet de Moradin. Les salamandres de feu ne doivent pas être autorisées à le conserver !

Si vous aimez le peuple nain, alors écoutez mon appel et rejoignez-moi pour rendre Hrakhamar à ses propriétaires légitimes. Je renoncerai à mes honoraires en échange de votre aide.
\end{quote}

Musharib proposait lui aussi de renoncer à son paiement, mais sa quête les entraînerait vers une forge occupée par des créatures de feu. Pour Morrÿs, l'appel d'un autre nain et la récupération d'une relique de Moradin ne manquaient cependant pas d'intérêt.

\subsection*{Qawasha et Kupalué}

\illustration[0.85\textwidth]{2026-09-22/prospectus/Qawasha et Kupalué.png}
{\textit{Prospectus de Qawasha et de son compagnon Kupalué.}}
{prospectus_qawasha}

\begin{quote}\itshape
Moi, Qawasha, je suis druide et membre de l'Enclave d'Émeraude. Je cherche à débarrasser Chult, ma terre natale, de la menace des morts-vivants qui menace d'éradiquer toute vie. Mon compagnon végépygmée connaît parfaitement la flore et la faune locales.

Chult est notre foyer, et nous pouvons vous guider comme personne d'autre.

Notre tarif est de cinq pièces d'or par jour, mais ce montant est négociable si vous acceptez de nous aider à détruire tous les morts-vivants que nous pourrions rencontrer. Peut-être pourrons-nous ensemble recevoir la bénédiction des esprits chwinga qui veillent sur cette terre sacrée.
\end{quote}

Destos connaissait la réputation de l'Enclave d'Émeraude. Ses membres cherchaient à préserver l'équilibre naturel, quitte à employer des méthodes sévères contre tout ce qui le menaçait. Leur hostilité envers les morts-vivants ne laissait aucun doute, et leur objectif pouvait rejoindre celui de l'expédition.

\subsection*{Salida}

\illustration[0.85\textwidth]{2026-09-22/prospectus/Salida.png}
{\textit{Prospectus de Salida.}}
{prospectus_salida}

\begin{quote}\itshape
J'ai mené de nombreuses expéditions à travers les terres sauvages de Chult. Je connais des sentiers secrets qui n'apparaissent sur aucune carte. Engagez les autres à vos risques et périls.

Azaka et Eku ont mené des expéditions dans la jungle qui ne sont jamais revenues. Faroul et Gondolo sont des charlatans, purement et simplement. Le vieux Hew souffre de la fièvre de la brume bleue, j'en suis certaine. Musharib est un nain albinos, et il ne faut pas lui faire confiance. Et j'ai entendu dire que River Mist et Flask of Wine doivent de l'argent au Zhentarim.

Cinq pièces d'or par jour, c'est un faible prix à payer pour vos vies. Dépensez votre or intelligemment !
\end{quote}

Obélik connaissait le Zhentarim sous son autre nom, le Réseau noir. Ses membres vendaient leurs services tant qu'ils pouvaient y gagner de l'or ou de l'influence et préféraient agir loin des regards. Cette réputation ne prouvait rien sur River Mist et Flask of Wine, mais elle n'encourageait guère à répondre à leur annonce clandestine.

Salida ne convainquit pas davantage Kael. Son prospectus consacrait plus de place à discréditer ses concurrents qu'à démontrer ses propres qualités. Peut-être disait-elle vrai, peut-être pas ; dans les deux cas, dénigrer les autres ne suffisait pas à faire d'elle la meilleure.

Les prospectus avaient réduit le nombre de possibilités sans offrir encore de réponse définitive. Plusieurs guides proposaient une réduction en échange d'une mission, mais chacune de ces dettes conduisait vers un nouveau danger. Il leur faudrait désormais rencontrer les candidats, comparer leurs itinéraires et choisir celui auquel ils accepteraient de confier leur vie.

La jungle pouvait attendre encore un peu.
